\documentclass[12pt]{article}
\usepackage{geometry}
\usepackage{amsmath}
\usepackage{amsfonts}
\usepackage{enumitem}
\usepackage{geometry}% 1-inch margins
\geometry{letterpaper, margin=1in}% 1-inch margins
\title{Extra Exercise \#2}
\author{Jacob C. Vaught}
\date{Due September 19th, 2023}

\begin{document}
\maketitle

\section*{Submission Guidelines}

This assignment requires a typed submission; handwritten responses will not be accepted. Make sure to \textbf{properly cite ALL the sources} from which you've gathered data, either as footnotes or in a references section at the end of your document.
\newline

Each problem in the assignment should be well-organized and consist of the following elements:

\begin{itemize}
  \item Problem Statement: Clearly outline the question or task you're attempting to solve.
  \item Equations Used: List all the relevant equations that will be employed in solving the problem.
  \item Detailed Solutions: Provide an in-depth, step-by-step solution to each problem, ensuring that all mathematical steps are clearly shown.
\end{itemize}

The quality of your writing is important. Treat this assignment as if it were a formal lab report. This includes correct grammar, coherent sentence structures, and professional formatting. To ensure the highest quality, it is recommended to utilize a large language model, such as ChatGPT, for proofreading your essays and solutions.
\newline

To use ChatGPT for proofreading, you can enter a query like "Rewrite this in passive, lab report format: [insert essay text here]" to receive improved sentence structures and grammar corrections.
\newline

Please submit your completed assignment on Blackboard as a PDF document. Failure to comply with these guidelines may result in a reduction of your grade.

\newpage
\section*{Problem 1: Vehicle Aerodynamics}
Calculate the drag force for multiple vehicles (minimum of three) from different categories, such as sedans, SUVs, electric cars, trucks, etc. Use these calculations to determine the change in miles per gallon (mpg) for each vehicle type. Assume that all other factors like air density, speed, and weight are constant. After conducting the calculations, write an essay comparing the aerodynamic efficiencies of these different vehicle categories. Discuss why electric cars generally have a lower Drag Coefficient (\(C_d\)) than vehicles with combustion engines.

\subsection*{Equations}
\begin{align*}
F_{\text{drag}} &= 0.5 \times C_d \times A \times \rho \times V^2 \\
\Delta \text{mpg} &= \frac{F_{\text{drag,old}}}{F_{\text{drag,new}}} \times \text{mpg}_{\text{old}}
\end{align*}

\newpage
\subsection*{Example Solution}
For illustration, let's look at two different kinds of vehicles—a sedan and an SUV:

\begin{enumerate}[label=\alph*.]
    \item Honda Civic (Sedan): \(C_d = 0.28\), \(A = 2.2 \, \text{m}^2\)
    \item Ford Expedition (SUV): \(C_d = 0.36\), \(A = 2.8 \, \text{m}^2\)
\end{enumerate}

We assume a constant air density of \(\rho = 1.225 \, \text{kg/m}^3\) and a speed of \(V = 30 \, \text{m/s}\).

First, we calculate the drag force for each using the formula \(F_{\text{drag}} = 0.5 \times C_d \times A \times \rho \times V^2\).

\begin{align*}
F_{\text{drag, Civic}} &= 0.5 \times 0.28 \times 2.2 \times 1.225 \times 30^2 = 252.3 \, \text{N} \\
F_{\text{drag, Expedition}} &= 0.5 \times 0.36 \times 2.8 \times 1.225 \times 30^2 = 467.4 \, \text{N}
\end{align*}

Subsequently, we'll use the \(\Delta \text{mpg}\) formula to ascertain the change in fuel efficiency for the Ford Expedition in comparison to the Honda Civic, which we'll assume has an original fuel efficiency of 25 mpg.

\[
\Delta \text{mpg} = \frac{252.3}{467.4} \times 25 = 13.5 \, \text{mpg}
\]
\newline
Essay:

In this example, the Ford Expedition, an SUV, would have significantly lower fuel efficiency when compared to the Honda Civic, a sedan. The Expedition's mpg would be approximately 13.5 compared to the Civic's 25 mpg. This dramatic difference can be attributed to their respective drag coefficients and frontal areas. 

Electric cars typically have lower drag coefficients as they are designed with a greater focus on aerodynamic efficiency to maximize battery range. On the other hand, vehicles like SUVs and trucks that are equipped with combustion engines often have higher drag coefficients, leading to lower fuel efficiency.




\newpage
\section*{Problem 2: Thermal Conductivity}

\subsection*{Problem Statement}
Calculate the thermal conductivity for three different building materials and a single-pane window measuring \(1 \, \text{m} \times 2 \, \text{m}\). Assume a wall area of at least \(10 \, \text{m} \times 10 \, \text{m}\), although you may consider a larger wall area if desired. Assume the interior temperature based on your home setting (convert to Celsius or Kelvin) and select a realistic external temperature. 
\newline

After calculations, compose an essay discussing your findings and any peculiar observations. Which building material was most effective for keeping heat out? Is there a reason it is/isn't used in home construction today? Was the thermal conductivity of the window a contributing factor to older homes lacking open glass designs before the advent of air conditioning? 

\subsection*{Equations}
\[
Q = \frac{k \times A \times \Delta T}{d}
\]
\newpage
\subsection*{Example Solution}

For this example, let's assume the following building materials:

\begin{enumerate}[label=\alph*.]
    \item Concrete: \( k = 1.7 \, \text{W/m}\cdot\text{K}\), \( A = 10 \, \text{m}^2 \), \( \Delta T = 20 \, \text{K} \), \( d = 0.2 \, \text{m} \)
    \item Wood: \( k = 0.16 \, \text{W/m}\cdot\text{K}\), \( A = 12 \, \text{m}^2 \), \( \Delta T = 20 \, \text{K} \), \( d = 0.05 \, \text{m} \)
    \item Brick: \( k = 0.6 \, \text{W/m}\cdot\text{K}\), \( A = 15 \, \text{m}^2 \), \( \Delta T = 20 \, \text{K} \), \( d = 0.1 \, \text{m} \)
    \item Window Glass: \( k = 1.05 \, \text{W/m}\cdot\text{K}\), \( A = 2 \, \text{m}^2 \), \( \Delta T = 20 \, \text{K} \), \( d = 0.005 \, \text{m} \)
\end{enumerate}

Calculating the thermal conductivity using the equation \(Q = \frac{k \times A \times \Delta T}{d}\):

\begin{align*}
Q_{\text{Concrete}} &= \frac{1.7 \times 10 \times 20}{0.2} = 1700 \, \text{W} \\
Q_{\text{Wood}} &= \frac{0.16 \times 12 \times 20}{0.05} = 768 \, \text{W} \\
Q_{\text{Brick}} &= \frac{0.6 \times 15 \times 20}{0.1} = 1800 \, \text{W} \\
Q_{\text{Window Glass}} &= \frac{1.05 \times 2 \times 20}{0.005} = 8400 \, \text{W}
\end{align*}
\newline
Essay:

Based on these calculations, we observe a striking difference in thermal conductivity across different materials. The most glaring contrast is with window glass, which allows for a heat flow of 8400 W. This suggests that window glass is highly inefficient in terms of thermal insulation when compared to materials like concrete, wood, or brick.

This inefficiency may offer an explanation as to why older homes, especially those built before air conditioning became commonplace, avoided large, open glass designs. The poor thermal insulation would make it challenging to maintain a comfortable internal temperature, especially in extreme weather conditions.


\newpage
\section*{Problem 3: Insulated Homes and Multi-pane Windows}
Investigate the thermal properties of a residential home by choosing two types of insulation materials, each with a thickness you deem standard (for example, 3.5 inches). Calculate their thermal conductivities. Then, integrate these insulations with the most thermally efficient building material you identified in Problem 2, to evaluate the overall thermal conductivity of the wall.
\newline

Next, turn your attention to windows. Using the material selected for the window in Problem 2, calculate the heat transfer properties of double- and triple-pane windows. Assume the second and third panes in these windows have the same material properties as the first.
\newline

As you perform these calculations, also consider the advancements in window technology. Double and triple pane windows typically use an inert gas. In your accompanying essay, explain the rationale behind this. Additionally, identify the type of inert gas most commonly used in modern multi-pane windows.
\newline

Conclude your essay by discussing the potential impact of multi-pane window technology on architectural design. Specifically, address whether the advent of multi-pane windows could explain the prevalence of larger window areas in modern buildings compared to their older counterparts.

\subsection*{Equations}
\[
Q = \frac{k \times A \times \Delta T}{d}
\]

\newpage
\subsection*{Example Solution}
\subsubsection*{Insulation Materials}

From Problem 2, we identified Brick as the most thermally efficient building material with a thermal conductivity of \( k = 0.6 \, \text{W/m}\cdot\text{K} \).

We will consider Cellulose as the insulation material for this example, with \( k = 0.035 \, \text{W/m}\cdot\text{K} \) and a thickness of 3.5 inches (\(d \approx 0.089 \, \text{m}\)). The Brick wall is assumed to have a thickness of 9 inches (\(d \approx 0.2286 \, \text{m}\)). The temperature difference \(\Delta T\) is given as \(15 \, \text{K}\).

Firstly, let's calculate the total thermal resistance (\( R_{\text{total}} \)) of the composite wall, which consists of both Cellulose and Brick, using the equation for thermal resistances in series:

\[
R_{\text{total}} = R_{\text{Brick}} + R_{\text{Cellulose}} = \frac{d_{\text{Brick}}}{k_{\text{Brick}}} + \frac{d_{\text{Cellulose}}}{k_{\text{Cellulose}}}
\]
\[
R_{\text{total}} = \frac{0.2286}{0.6} + \frac{0.089}{0.035} \approx 0.381 + 2.543 = 2.924 \, \text{m}\cdot\text{K/W}
\]
Now, the heat transfer (\( Q \)) through the composite wall can be calculated as follows:

\[
Q_{\text{Composite}} = \frac{\Delta T}{R_{\text{total}}} = \frac{15}{2.924} \approx 5.13 \, \text{W}
\]
To understand the efficacy of the insulation, we should also calculate the heat transfer (\( Q \)) for just the Brick wall:

\[
Q_{\text{Brick}} = \frac{0.6 \times 1 \times 15}{0.2286} \approx 39.43 \, \text{W}
\]

By using Cellulose as insulation, the heat transfer rate is reduced from approximately \( 39.43 \, \text{W} \) for the Brick wall alone to \( 5.13 \, \text{W} \) for the composite wall. This illustrates the significant improvement in thermal efficiency due to the insulation.


\subsubsection*{Multi-pane Windows}

In Problem 2, we identified Window Glass with a thermal conductivity of \(k = 1.05 \, \text{W/m} \cdot \text{K}\). We'll use this for double- and triple-pane windows, each pane having a thickness of 0.005 m. Let's also assume an inert gas with \(k = 0.016 \, \text{W/m} \cdot \text{K}\) filling the gap with the same thickness.

For Double-pane:
\[
Q_{\text{Double-pane}} = \frac{15}{2 \times \frac{0.005}{1.05} + \frac{0.005}{0.016}} \approx 1576.7 \, \text{W}
\]

For Triple-pane:
\[
Q_{\text{Triple-pane}} = \frac{15}{3 \times \frac{0.005}{1.05} + 2 \times \frac{0.005}{0.016}} \approx 1050.0 \, \text{W}
\]

\subsubsection*{Essay}

In modern times, XXXX is the most commonly used inert gas for this purpose.

Regarding architectural design, the invention of multi-pane windows with inert gases has indeed allowed for larger windows to be installed without as great a loss in thermal efficiency as single-pane windows. This could explain why modern buildings tend to feature larger windows than their older counterparts.

\end{document}